\chapter{Errata in the Official Solutions}
\label{ch:errata}

During transcription every solution was checked digit by digit against
the official papers. The list below collects the places where the
\emph{official} text itself is wrong or internally inconsistent ---
sign slips, inverted relations, mismatched values between a derivation
and its numerical substitution, misprints, and stray editorial leftovers.
In the body of this book each is transcribed \emph{as printed} and
marked with a comment in the source; this appendix is the reader-facing
summary. Trivial spelling slips are included only where they could
confuse (wrong symbol names, wrong units).

\section*{Chapter 1 --- Positional Astronomy \& Time}

\noindent\textbf{Solution 1.11} (IOAA 2012 TH S13): presumably 2*pi*t/T is meant
\smallskip

\noindent\textbf{Solution 1.15} (IOAA 2013 TH S11): stray "(P)" in source
\smallskip

\noindent\textbf{Solution 1.17} (IOAA 2013 TH S9): source gives the declination as 18$^\circ$05' above but uses 18$^\circ$07' here
\smallskip

\noindent\textbf{Solution 1.18} (IOAA 2014 TH L1): the solution switches from d\_ES to d\_PS for the Earth-Sun distance
\smallskip

\noindent\textbf{Solution 1.18} (IOAA 2014 TH L1): the problem gives phi = 45 degrees, but the solution evaluates with 45deg21'
\smallskip

\noindent\textbf{Solution 1.18} (IOAA 2014 TH L1): source writes delta\_min (not delta\_0,min) in the middle term
\smallskip

\noindent\textbf{Solution 1.19} (IOAA 2014 TH P13): source has stray "(G\_1;g\_1)" here and below, and never names the second triangle (presumably sigma\_1 B sigma\_2)
\smallskip

\noindent\textbf{Solution 1.25} (IOAA 2018 TH T10): source says "3617 BCE"; 363700 years ago would be \textasciitilde{}361700 BCE
\smallskip

\noindent\textbf{Solution 1.25} (IOAA 2018 TH T10): source labels the third row with Vega's components again; the numbers below are Altair's
\smallskip

\noindent\textbf{Solution 1.29} (IOAA 2018 TH T5): the official solution says 18 h; at the September equinox the Sun's R.A. is 12 h.
\smallskip

\noindent\textbf{Solution 1.33} (IOAA 2020 TH 7): the printed solution has sin(A'\_i - A) in the two left-hand terms; mathematically it should be sin(A'\_i - A\_i).
\smallskip

\noindent\textbf{Solution 1.33} (IOAA 2020 TH 7): the official solution writes -13.4 deg in two places here and below where the altitude is a = -13.96 deg.
\smallskip

\noindent\textbf{Solution 1.34} (IOAA 2021 TH Q8): the official solution writes d\_{1-2} twice in the denominator; it should be 2 d\_{1-2} d\_{1-3}.
\smallskip

\noindent\textbf{Solution 1.35} (IOAA 2022 TH 7): the official solution writes 9900 here; from the previous line it should be 9990 km.
\smallskip

\noindent\textbf{Solution 1.37} (IOAA 2023 TH Q10): the official solution writes 4416 here; the relative speed derived above is 4412 km/h.
\smallskip

\noindent\textbf{Solution 1.61} (IOAA 2015 OBS-SKY 1): M20 is listed twice in the key table
\smallskip

\noindent\textbf{Solution 1.61} (IOAA 2015 OBS-SKY 1): M24's listed magnitude and size as in the key
\smallskip

\noindent\textbf{Solution 1.72} (IOAA 2022 OBS-SKY O5): ``lense'' in the original
\smallskip

\noindent\textbf{Solution 1.74} (IOAA 2022 OBS-SKY O7): "Vulpecua" for Vulpecula in the source
\smallskip

\noindent\textbf{Solution 1.74} (IOAA 2022 OBS-SKY O7): "Aquatius" for Aquarius in the source
\smallskip

\noindent\textbf{Solution 1.78} (IOAA 2024 OBS-SKY 2): "Celeste Chart" in the source
\smallskip

\noindent\textbf{Solution 1.83} (IOAA 2011 OBS-TEL 1): "Delfinus" for Delphinus in the source
\smallskip

\noindent\textbf{Solution 1.83} (IOAA 2011 OBS-TEL 1): the source says "finder scope" here although Drawing 2 is the view through the larger telescope
\smallskip

\section*{Chapter 2 --- Celestial Mechanics \& Gravitation}

\noindent\textbf{Solution 2.5} (IOAA 2009 TH P12): source uses 3.19e-10 here although $\dot{r}$ was 3.21e-10 above
\smallskip

\noindent\textbf{Solution 2.10} (IOAA 2010 TH L16): source prints $\sqrt{2k_{sun}r_M}$; dimensionally it should be $\sqrt{2k_{sun}/r_M}$
\smallskip

\noindent\textbf{Solution 2.12} (IOAA 2011 TH S1): source prints a\^{}{2/3}; Kepler's third law gives T = a\^{}{3/2}
\smallskip

\noindent\textbf{Solution 2.14} (IOAA 2011 TH S4): source omits the cosine; the geometry (and the numerical value below) give cos(alpha) = R\_M/R\_P
\smallskip

\noindent\textbf{Solution 2.14} (IOAA 2011 TH S4): source prints sqrt(GM\_M / R\_P\^{}{3/2}); dimensionally it should be sqrt(GM\_M / R\_P\^{}3)
\smallskip

\noindent\textbf{Solution 2.15} (IOAA 2011 TH S6): the minus sign of the exponent is illegible in the scan
\smallskip

\noindent\textbf{Solution 2.15} (IOAA 2011 TH S6): change of omega, not epsilon
\smallskip

\noindent\textbf{Solution 2.18} (IOAA 2013 TH L2): "non-retrogate" in source
\smallskip

\noindent\textbf{Solution 2.19} (IOAA 2014 TH L2): "satelite"
\smallskip

\noindent\textbf{Solution 2.19} (IOAA 2014 TH L2): "suspention"
\smallskip

\noindent\textbf{Solution 2.19} (IOAA 2014 TH L2): "is negligible results"
\smallskip

\noindent\textbf{Solution 2.19} (IOAA 2014 TH L2): "solutions ... is"
\smallskip

\noindent\textbf{Solution 2.19} (IOAA 2014 TH L2): Romanian "stable equilibrium"
\smallskip

\noindent\textbf{Solution 2.19} (IOAA 2014 TH L2): Romanian "unstable equilibrium"
\smallskip

\noindent\textbf{Solution 2.19} (IOAA 2014 TH L2): "Coresponding"
\smallskip

\noindent\textbf{Solution 2.19} (IOAA 2014 TH L2): "circulae"
\smallskip

\noindent\textbf{Solution 2.19} (IOAA 2014 TH L2): "The potentialenergy v" as printed
\smallskip

\noindent\textbf{Solution 2.19} (IOAA 2014 TH L2): identical expression repeated on both sides
\smallskip

\noindent\textbf{Solution 2.19} (IOAA 2014 TH L2): the official solution switches to Romanian here
\smallskip

\noindent\textbf{Solution 2.19} (IOAA 2014 TH L2): F for F\_t inside the bracket
\smallskip

\noindent\textbf{Solution 2.19} (IOAA 2014 TH L2): part d) of the official solution is in Romanian
\smallskip

\noindent\textbf{Solution 2.20} (IOAA 2014 TH L3): "Eest"
\smallskip

\noindent\textbf{Solution 2.20} (IOAA 2014 TH L3): statement gives 44 deg North
\smallskip

\noindent\textbf{Solution 2.20} (IOAA 2014 TH L3): "EEst"
\smallskip

\noindent\textbf{Solution 2.20} (IOAA 2014 TH L3): "schetch"
\smallskip

\noindent\textbf{Solution 2.20} (IOAA 2014 TH L3): "elypse"
\smallskip

\noindent\textbf{Solution 2.20} (IOAA 2014 TH L3): Romanian "it is known that"
\smallskip

\noindent\textbf{Solution 2.20} (IOAA 2014 TH L3): Romanian "it results"
\smallskip

\noindent\textbf{Solution 2.20} (IOAA 2014 TH L3): missing factor 3 in the middle member
\smallskip

\noindent\textbf{Solution 2.20} (IOAA 2014 TH L3): Romanian "equilateral triangle"
\smallskip

\noindent\textbf{Solution 2.20} (IOAA 2014 TH L3): "Accordin to Kepplers"
\smallskip

\noindent\textbf{Solution 2.20} (IOAA 2014 TH L3): "Daca For", mixed Romanian/English
\smallskip

\noindent\textbf{Solution 2.20} (IOAA 2014 TH L3): "bellow", "enegy"
\smallskip

\noindent\textbf{Solution 2.20} (IOAA 2014 TH L3): the remainder of f) is largely in Romanian
\smallskip

\noindent\textbf{Solution 2.20} (IOAA 2014 TH L3): garbled subscript "Soa,,re" as printed
\smallskip

\noindent\textbf{Solution 2.20} (IOAA 2014 TH L3): "Symilarly"
\smallskip

\noindent\textbf{Solution 2.21} (IOAA 2014 TH P1): the figures are actually labelled "Figure 1a" and "Figure 1b" in the source
\smallskip

\noindent\textbf{Solution 2.21} (IOAA 2014 TH P1): decimal comma and capital "Km" as in source
\smallskip

\noindent\textbf{Solution 2.22} (IOAA 2014 TH P2): "decrees" (decrease) in source
\smallskip

\noindent\textbf{Solution 2.22} (IOAA 2014 TH P2): source writes r\_{0,max} here; should be r\_{0,min}
\smallskip

\noindent\textbf{Solution 2.22} (IOAA 2014 TH P2): "Keppler's" in source
\smallskip

\noindent\textbf{Solution 2.22} (IOAA 2014 TH P2): Romanian "este" (is) left in source
\smallskip

\noindent\textbf{Solution 2.22} (IOAA 2014 TH P2): "Erath" in source
\smallskip

\noindent\textbf{Solution 2.23} (IOAA 2014 TH P5): relation (5) is a verbatim repeat of relation (4) in the source
\smallskip

\noindent\textbf{Solution 2.23} (IOAA 2014 TH P5): "gass" in source
\smallskip

\noindent\textbf{Solution 2.23} (IOAA 2014 TH P5): "accelearated" in source
\smallskip

\noindent\textbf{Solution 2.23} (IOAA 2014 TH P5): "in at the shuttle", "to much" in source
\smallskip

\noindent\textbf{Solution 2.24} (IOAA 2015 TH L1): "when it closest"
\smallskip

\noindent\textbf{Solution 2.24} (IOAA 2015 TH L1): "R(pi - alpha\_0 alpha\_0)" as printed
\smallskip

\noindent\textbf{Solution 2.24} (IOAA 2015 TH L1): stray "= 900" inside the bracket as printed
\smallskip

\noindent\textbf{Solution 2.24} (IOAA 2015 TH L1): "in terms new variables"
\smallskip

\noindent\textbf{Solution 2.24} (IOAA 2015 TH L1): exponent missing on (y*) as printed
\smallskip

\noindent\textbf{Solution 2.24} (IOAA 2015 TH L1): right-hand side left blank as printed
\smallskip

\noindent\textbf{Solution 2.24} (IOAA 2015 TH L1): "the the"
\smallskip

\noindent\textbf{Solution 2.24} (IOAA 2015 TH L1): both the second and third versions are titled "Cartesian version"
\smallskip

\noindent\textbf{Solution 2.24} (IOAA 2015 TH L1): "semi major" twice as printed
\smallskip

\noindent\textbf{Solution 2.24} (IOAA 2015 TH L1): heading repeated from (a) as printed
\smallskip

\noindent\textbf{Solution 2.24} (IOAA 2015 TH L1): "characterized"
\smallskip

\noindent\textbf{Solution 2.24} (IOAA 2015 TH L1): "it is at necessary"
\smallskip

\noindent\textbf{Solution 2.24} (IOAA 2015 TH L1): trailing "=" as printed
\smallskip

\noindent\textbf{Solution 2.27} (IOAA 2015 TH S13): arguments written with R\_Io, the fractions below use R\_Io/2
\smallskip

\noindent\textbf{Solution 2.27} (IOAA 2015 TH S13): exact denominator would be (1-(R\_Io/2r)\^{}2)\^{}2
\smallskip

\noindent\textbf{Solution 2.27} (IOAA 2015 TH S13): "If assume"
\smallskip

\noindent\textbf{Solution 2.27} (IOAA 2015 TH S13): 6R\_Io/8r in source; exact would be 3R\_Io/8r
\smallskip

\noindent\textbf{Solution 2.28} (IOAA 2015 TH S2): "satelite" in source
\smallskip

\noindent\textbf{Solution 2.28} (IOAA 2015 TH S2): the source labels this ratio (and the next) rho\_E/rho\_P; the value 0.238 is actually rho\_P/rho\_E, as the Second Alternate Version confirms
\smallskip

\noindent\textbf{Solution 2.28} (IOAA 2015 TH S2): "Plant" in source
\smallskip

\noindent\textbf{Solution 2.28} (IOAA 2015 TH S2): the source writes m\_E/m\_P for what is actually m\_P/m\_E
\smallskip

\noindent\textbf{Solution 2.29} (IOAA 2015 TH S5): sign convention with perihelion at theta = pi
\smallskip

\noindent\textbf{Solution 2.29} (IOAA 2015 TH S5): the 0.5 factor is dropped on the right-hand sides
\smallskip

\noindent\textbf{Solution 2.31} (IOAA 2016 TH T9): the source prints the denominator as "(3.0275 + 6.635 x 10\^{}6)"; it should read (3.0275 x 10\^{}7 + 6.635 x 10\^{}6).
\smallskip

\noindent\textbf{Solution 2.32} (IOAA 2019 TH 9): the source writes "xR\_Moon" although 384 400 km is the Earth--Moon distance R, not the radius of the Moon.
\smallskip

\noindent\textbf{Solution 2.35} (IOAA 2021 TH Q12): the source prints "Gm" here rather than "GM".
\smallskip

\noindent\textbf{Solution 2.36} (IOAA 2021 TH Q14): "[2 point]"
\smallskip

\noindent\textbf{Solution 2.36} (IOAA 2021 TH Q14): "[1 points]"
\smallskip

\noindent\textbf{Solution 2.36} (IOAA 2021 TH Q14): "[2 point]"
\smallskip

\noindent\textbf{Solution 2.36} (IOAA 2021 TH Q14): "[1 points]"
\smallskip

\noindent\textbf{Solution 2.36} (IOAA 2021 TH Q14): "[1 points]"
\smallskip

\noindent\textbf{Solution 2.36} (IOAA 2021 TH Q14): "[3 point]"
\smallskip

\noindent\textbf{Solution 2.37} (IOAA 2021 TH Q3): "braking"
\smallskip

\noindent\textbf{Solution 2.39} (IOAA 2022 TH 12): "are opposing positions"
\smallskip

\noindent\textbf{Solution 2.39} (IOAA 2022 TH 12): "way straightforward"
\smallskip

\noindent\textbf{Solution 2.39} (IOAA 2022 TH 12): upright r1, r2 in source
\smallskip

\noindent\textbf{Solution 2.39} (IOAA 2022 TH 12): "the into given"
\smallskip

\noindent\textbf{Solution 2.40} (IOAA 2022 TH 8): "at at"
\smallskip

\noindent\textbf{Solution 2.41} (IOAA 2023 TH Q12): "the mass of the spacecraft mass"
\smallskip

\noindent\textbf{Solution 2.41} (IOAA 2023 TH Q12): 0.17662 for 0.17622
\smallskip

\noindent\textbf{Solution 2.42} (IOAA 2024 TH T11): "(improve figure later)" left in the official solutions
\smallskip

\noindent\textbf{Solution 2.42} (IOAA 2024 TH T11): "as a function of the polar radius"
\smallskip

\noindent\textbf{Solution 2.43} (IOAA 2024 TH T3): "focii", "apapsis"
\smallskip

\noindent\textbf{Solution 2.44} (IOAA 2020 DA 2): "a maximum and maximum"
\smallskip

\noindent\textbf{Solution 2.44} (IOAA 2020 DA 2): 2a = 2.98 au (cf. 2a\_P = 5.965 au above)
\smallskip

\noindent\textbf{Solution 2.44} (IOAA 2020 DA 2): units given as yr in source, should be au
\smallskip

\section*{Chapter 3 --- Solar System \& Planetary Science}

\noindent\textbf{Solution 3.10} (IOAA 2010 TH S14): the factor cos(epsilon) appears in the source without epsilon being defined anywhere in the solution.
\smallskip

\noindent\textbf{Solution 3.11} (IOAA 2010 TH S15): the source defines "t2" as the duration of the total eclipse although the formula above uses t1.
\smallskip

\noindent\textbf{Solution 3.12} (IOAA 2011 TH S10): "should be pass", "larger then" as in the original.
\smallskip

\noindent\textbf{Solution 3.16} (IOAA 2014 TH P14): subscripts C,S / C,O here vs B,S / B,O in the text below, as in the original
\smallskip

\noindent\textbf{Solution 3.16} (IOAA 2014 TH P14): "unde" (Romanian for "where") as in the original
\smallskip

\noindent\textbf{Solution 3.16} (IOAA 2014 TH P14): "bellow" as in the original
\smallskip

\noindent\textbf{Solution 3.16} (IOAA 2014 TH P14): "particularaly" as in the original
\smallskip

\noindent\textbf{Solution 3.16} (IOAA 2014 TH P14): "SU" as in the original
\smallskip

\noindent\textbf{Solution 3.16} (IOAA 2014 TH P14): 65536/0.2 = 327680, source prints 491520
\smallskip

\noindent\textbf{Solution 3.17} (IOAA 2015 TH S3): (1 + 3.076\^{}2 - 2.156\^{}2)/(2 x 3.076) = 0.945; source prints 0.995, while the numbers below are consistent with 0.945
\smallskip

\noindent\textbf{Solution 3.17} (IOAA 2015 TH S3): "t\_occil" as in the original; 13984.16 x 6.201 = 86715.6, source prints 86701.8
\smallskip

\noindent\textbf{Solution 3.18} (IOAA 2015 TH S9): "e-5" calculator notation and unrounded digit strings as in the original
\smallskip

\noindent\textbf{Solution 3.18} (IOAA 2015 TH S9): the solid angle uses the full 66 arcsec angular diameter in place of the 33 arcsec radius, as in the original
\smallskip

\noindent\textbf{Solution 3.20} (IOAA 2017 TH T8): "will locate elongation" as in the original
\smallskip

\noindent\textbf{Solution 3.20} (IOAA 2017 TH T8): 102 x 242.7 = 24755.4, source prints 24758.33
\smallskip

\noindent\textbf{Solution 3.27} (IOAA 2023 TH Q7): "the students does"
\smallskip

\noindent\textbf{Solution 3.28} (IOAA 2024 TH T10): units given as m in the source
\smallskip

\noindent\textbf{Solution 3.28} (IOAA 2024 TH T10): 77.59 here, though the middle step and the modulus below use 77.76
\smallskip

\noindent\textbf{Solution 3.30} (IOAA 2007 DA D3): elongation given as 181.20 degrees in the source
\smallskip

\noindent\textbf{Solution 3.31} (IOAA 2008 DA D3): 0.0013487 listed, though the formula below gives 0.0013479
\smallskip

\noindent\textbf{Solution 3.31} (IOAA 2008 DA D3): 0.0060566 listed, though the formula below gives 0.0060532
\smallskip

\noindent\textbf{Solution 3.31} (IOAA 2008 DA D3): 0.095 and 0.435 both appear in the source
\smallskip

\noindent\textbf{Solution 3.34} (IOAA 2014 DA D2): "hav"
\smallskip

\noindent\textbf{Solution 3.34} (IOAA 2014 DA D2): "of de gas"
\smallskip

\noindent\textbf{Solution 3.34} (IOAA 2014 DA D2): "is spend"
\smallskip

\noindent\textbf{Solution 3.34} (IOAA 2014 DA D2): "rapport"
\smallskip

\noindent\textbf{Solution 3.34} (IOAA 2014 DA D2): "grah"
\smallskip

\noindent\textbf{Solution 3.34} (IOAA 2014 DA D2): Romanian "?i" (and) left in the original
\smallskip

\noindent\textbf{Solution 3.34} (IOAA 2014 DA D2): "plot de graph"
\smallskip

\noindent\textbf{Solution 3.34} (IOAA 2014 DA D2): the table prints 593,3 K although the computation above gives 593,7 K.
\smallskip

\noindent\textbf{Solution 3.35} (IOAA 2016 DA D2): the boxed answer prints a doubled equals sign.
\smallskip

\noindent\textbf{Solution 3.37} (IOAA 2025 OBS-PLAN OP02): "upto"
\smallskip

\section*{Chapter 4 --- The Sun \& Heliophysics}

\noindent\textbf{Solution 4.1} (IOAA 2008 TH S2): 32 x 60"/1.86" = 1032 pixels; the official solution states 1034 pixels
\smallskip

\noindent\textbf{Solution 4.5} (IOAA 2018 TH T7): "larger that" as in the original
\smallskip

\noindent\textbf{Solution 4.6} (IOAA 2019 TH 7): "when the spot directed towards" as in the original
\smallskip

\noindent\textbf{Solution 4.7} (IOAA 2021 TH Q5): the official solution writes "At z = H, B = B\_0 e\^{}{-z/150}"; with H = 150 km the exponent should read -z/(2 x 150) = -z/300, as used below.
\smallskip

\noindent\textbf{Solution 4.7} (IOAA 2021 TH Q5): ln 10 = 2.303; the official solution uses 2.301
\smallskip

\noindent\textbf{Solution 4.9} (IOAA 2025 DA D02): the official solution writes v\_{1,m} for CME-2 as well (should be v\_{2,m})
\smallskip

\noindent\textbf{Solution 4.9} (IOAA 2025 DA D02): "error is arrival time" as in the original
\smallskip

\noindent\textbf{Solution 4.9} (IOAA 2025 DA D02): "model start" as in the original
\smallskip

\noindent\textbf{Solution 4.10} (IOAA 2021 OBS-SKY 1): the official solution lists 2021/05/29 23:42:08 UT here, while its summary table gives 2021/05/28 23:42:08 UT for interval 4
\smallskip

\noindent\textbf{Solution 4.10} (IOAA 2021 OBS-SKY 1): the summary table gives 2021/05/29 05:36:08 UT for interval 7, but 05:00:08 UT is consistent with the quoted 401 minutes from onset
\smallskip

\noindent\textbf{Solution 4.10} (IOAA 2021 OBS-SKY 1): the "Time interval" column of the official table is not consistent with the differences of the successive timestamps (rows 4--7)
\smallskip

\noindent\textbf{Solution 4.10} (IOAA 2021 OBS-SKY 1): "difuminate" as in the original
\smallskip

\noindent\textbf{Solution 4.10} (IOAA 2021 OBS-SKY 1): the official solution uses 511 m/s (instead of the CME speed of 511 km/s); the quoted eV values follow from the m/s figure.
\smallskip

\section*{Chapter 5 --- Optics, Telescopes \& Detectors}

\noindent\textbf{Solution 5.2} (IOAA 2008 TH S10): the printed denominator is sin(6'.1020); numerically the quoted result 0.9388 corresponds to dividing by sin(gamma) = sin(6'.3424).
\smallskip

\noindent\textbf{Solution 5.6} (IOAA 2011 TH S5): the source writes sin(lambda) for the sine of the resolution angle
\smallskip

\noindent\textbf{Solution 5.17} (IOAA 2019 TH 14): the source prints omega\_r = v\_r / r; for the angular speed seen from the observing site directly below, the distance should be the altitude h.
\smallskip

\noindent\textbf{Solution 5.20} (IOAA 2021 TH Q13): "driven" (presumably "excited") as printed in the source.
\smallskip

\noindent\textbf{Solution 5.21} (IOAA 2021 TH Q4): the source spells "Jules" for Joules.
\smallskip

\noindent\textbf{Solution 5.21} (IOAA 2021 TH Q4): "temperatura" appears in the source.
\smallskip

\noindent\textbf{Solution 5.24} (IOAA 2024 TH T6): the source prints 12.8 here although V\_sky = 12.9 mag above.
\smallskip

\noindent\textbf{Solution 5.25} (IOAA 2025 TH T01): "can localized" as printed in the source.
\smallskip

\noindent\textbf{Solution 5.26} (IOAA 2025 TH T09): the source says "The radius of the image is therefore" although d\_image is the diameter asked for and the boxed value is the diameter.
\smallskip

\noindent\textbf{Solution 5.29} (IOAA 2013 DA D1): the source numbers this step "4." again, duplicating the previous item's number
\smallskip

\noindent\textbf{Solution 5.29} (IOAA 2013 DA D1): the source sentence "(If errors are included" is left unfinished
\smallskip

\noindent\textbf{Solution 5.30} (IOAA 2024 DA D1): the solution table gives Dec 0.079747 for ID2 = 11, whereas Table 2 of the problem lists 0.079474
\smallskip

\noindent\textbf{Solution 5.48} (IOAA 2022 OBS-TEL O1): the official solution table spells the name "Antilicus"; the question paper has "Antilicus" as well (the star is usually called Kornephoros)
\smallskip

\noindent\textbf{Solution 5.50} (IOAA 2022 OBS-TEL O3): the official solution writes "f" for the telescope focal length in the quoted range although the formula calls it F
\smallskip

\noindent\textbf{Solution 5.52} (IOAA 2023 OBS-TEL 1): the official solution uses "Delta x" in the bands for the error in mid-time although the heading defines it as Delta t
\smallskip

\noindent\textbf{Solution 5.53} (IOAA 2025 OBS-TEL OT01): the official marking table says "objective" for the eyepiece
\smallskip

\section*{Chapter 6 --- Radiation, Magnitudes \& Spectra}

\noindent\textbf{Problem 6.16} (IOAA 2014 TH P7): the exponent below is printed as hc/k$\lambda$ c in the original paper (physically it should be hc/$\lambda$ kT); transcribed as printed.
\smallskip

\noindent\textbf{Problem 6.17} (IOAA 2014 TH P8): the original paper prints the exponent as hc/k$\lambda$ c (presumably a typo for hc/k$\lambda$ T; cf. the closing condition below).
\smallskip

\noindent\textbf{Problem 6.29} (IOAA 2014 DA D3): the symbol in the original prints as a stray grave accent; it is most plausibly a script ell.
\smallskip

\noindent\textbf{Solution 6.6} (IOAA 2008 TH L2): "det" (detik, Indonesian for second) in the source
\smallskip

\noindent\textbf{Solution 6.9} (IOAA 2011 TH S8): "form" for "from"
\smallskip

\noindent\textbf{Solution 6.11} (IOAA 2012 TH L1): "astronomer 2" in the source; this is the flux seen by the observer at the centre (the biologist)
\smallskip

\noindent\textbf{Solution 6.12} (IOAA 2012 TH S9): the source says "U - V = 0.300"; the problem gives B - V = 0.300
\smallskip

\noindent\textbf{Solution 6.12} (IOAA 2012 TH S9): the values 0.038 and 0.365 on these two lines are swapped in the source; with the given data (U-V)\_1 = 0.365 and (U-B)\_1 = 0.038
\smallskip

\noindent\textbf{Solution 6.13} (IOAA 2013 TH S13): source has the Greek word for "or" here
\smallskip

\noindent\textbf{Solution 6.15} (IOAA 2014 TH P11): "hole" for "whole"
\smallskip

\noindent\textbf{Solution 6.15} (IOAA 2014 TH P11): the source labels this equation (2) although (2) was already used above
\smallskip

\noindent\textbf{Solution 6.16} (IOAA 2014 TH P7): the source prints k lambda\_2 T in both numerator and denominator; the denominator should be k lambda\_1 T
\smallskip

\noindent\textbf{Solution 6.18} (IOAA 2014 TH P9): the source writes "planet" subscripts although the problem concerns the Sun and a dust particle
\smallskip

\noindent\textbf{Solution 6.18} (IOAA 2014 TH P9): "cL\_S\^{}2" in the source; from relation (1) and the following lines the denominator should be c D\_S\^{}2
\smallskip

\noindent\textbf{Solution 6.18} (IOAA 2014 TH P9): the final formula omits the factor 1/c which follows from the preceding line; dimensionally d = (3/2) sigma T\_S\^{}4 R\_S\^{}2 / (c rho K M\_S)
\smallskip

\noindent\textbf{Solution 6.21} (IOAA 2017 TH T4): the numerator's units are printed "km\^{}-1 Mpc\^{}-1" in the source
\smallskip

\noindent\textbf{Solution 6.23} (IOAA 2018 TH T6): the source prints a factor pi\^{}2 here and in the next line; the collecting area of a dish of effective diameter d\_c is pi d\_c\^{}2/4
\smallskip

\noindent\textbf{Solution 6.24} (IOAA 2019 TH 5): the source says "by Delta t"; the temperature rise is Delta T
\smallskip

\noindent\textbf{Solution 6.25} (IOAA 2022 TH 11): "We know the that asteroid" in source
\smallskip

\noindent\textbf{Solution 6.27} (IOAA 2025 TH T04): the source substitutes 0.69 here although E(B-V) = 0.73 was obtained above (full credit was given for E(B-V) between 0.69 and 0.73)
\smallskip

\noindent\textbf{Solution 6.28} (IOAA 2013 DA D3): dash missing
\smallskip

\noindent\textbf{Solution 6.29} (IOAA 2014 DA D3): "Sirrius" in source
\smallskip

\noindent\textbf{Solution 6.29} (IOAA 2014 DA D3): "luminosties"
\smallskip

\noindent\textbf{Solution 6.29} (IOAA 2014 DA D3): "Symilarlyfor"
\smallskip

\noindent\textbf{Solution 6.29} (IOAA 2014 DA D3): 0,134 as printed; the data of the problem (m = 3,3 at 58 ly) give L\_4 of order 13 L\_S
\smallskip

\noindent\textbf{Solution 6.29} (IOAA 2014 DA D3): source prints "between dintre Big Dipper"
\smallskip

\noindent\textbf{Solution 6.29} (IOAA 2014 DA D3): "dreptunghic" (Romanian: right-angled)
\smallskip

\noindent\textbf{Solution 6.29} (IOAA 2014 DA D3): 55,5' for seconds as printed
\smallskip

\noindent\textbf{Solution 6.29} (IOAA 2014 DA D3): as printed
\smallskip

\noindent\textbf{Solution 6.29} (IOAA 2014 DA D3): "Eaarth" in source
\smallskip

\noindent\textbf{Solution 6.29} (IOAA 2014 DA D3): as printed
\smallskip

\section*{Chapter 7 --- Stellar Astrophysics}

\noindent\textbf{Solution 7.9} (IOAA 2013 TH S6): 19.15 vs 19.05 in source
\smallskip

\noindent\textbf{Solution 7.10} (IOAA 2013 TH S7): source says "barycentric" for bolometric
\smallskip

\noindent\textbf{Solution 7.11} (IOAA 2014 TH P10): source writes aT here and below (should be aT\^{}4)
\smallskip

\noindent\textbf{Solution 7.13} (IOAA 2015 TH L3): source flips to (T\_i - T\_f) in the second line
\smallskip

\section*{Chapter 8 --- Binary \& Variable Stars}

\noindent\textbf{Solution 8.6} (IOAA 2013 TH S4): the period is read off the light curve, i.e. Figure 3 of the statement
\smallskip

\noindent\textbf{Solution 8.6} (IOAA 2013 TH S4): the period-luminosity relation is Figure 2 of the statement
\smallskip

\noindent\textbf{Solution 8.6} (IOAA 2013 TH S4): the extinction term should be subtracted (m - M - A = -5 + 5 log r), which gives log r = 4.68 and r \textasciitilde{} 48 kpc.
\smallskip

\noindent\textbf{Solution 8.6} (IOAA 2013 TH S4): 10\^{}{4.78} = 60300 pc = 60.3 kpc; neither printed value matches, and 53000 pc is inconsistent with 64.5 kpc.
\smallskip

\noindent\textbf{Solution 8.7} (IOAA 2014 TH P15): the source prints the exponent 3/2; consistency with the earlier R\^{}2 = (KM/4pi\^{}2)\^{}{2/3} P\^{}{4/3} requires 2/3.
\smallskip

\noindent\textbf{Solution 8.12} (IOAA 2019 TH 8): expanding -2.5 log I gives +2.5 log e C\_b/(lambda T); the sign slip is absorbed by the absolute value in (8.3).
\smallskip

\noindent\textbf{Solution 8.15} (IOAA 2022 TH 6): the source writes M\_B again; this last line is clearly M\_A (= 40.7 - 11.4 = 29.3 solar masses).
\smallskip

\noindent\textbf{Solution 8.16} (IOAA 2023 TH Q9): the source writes "m\_{1,2} = m\_1 + m\_2"; m\_{1,2} is really the combined magnitude of both stars, not a literal sum of magnitudes.
\smallskip

\noindent\textbf{Solution 8.18} (IOAA 2008 DA D2): the exponent shows 5.4949 where Step 3 gave M\_bolB = 5.4849; the printed result 0.8763 follows from 5.4849.
\smallskip

\noindent\textbf{Solution 8.18} (IOAA 2008 DA D2): the source writes 1.10 here where every other iteration uses m\_bolB = 1.09.
\smallskip

\noindent\textbf{Solution 8.20} (IOAA 2011 DA D1): so printed in the source; the first term should read 2454092.4111, giving JD 2455806.5733.
\smallskip

\noindent\textbf{Solution 8.22} (IOAA 2015 DA D1): the source uses F\_sun = 2.913e-10 here although the value computed above was 3.31e-10 W/m\^{}2.
\smallskip

\section*{Chapter 9 --- Exoplanets}

\noindent\textbf{Solution 9.2} (IOAA 2011 TH L1): source misprint for "expression"
\smallskip

\noindent\textbf{Solution 9.2} (IOAA 2011 TH L1): source says "km/h"; the value is in km/s
\smallskip

\noindent\textbf{Solution 9.2} (IOAA 2011 TH L1): sentence truncated in source
\smallskip

\noindent\textbf{Solution 9.3} (IOAA 2012 TH S11): source has "lieexactly"
\smallskip

\noindent\textbf{Solution 9.7} (IOAA 2021 TH Q10): Spanish "days" left in the source
\smallskip

\noindent\textbf{Solution 9.8} (IOAA 2021 TH Q6): source wording
\smallskip

\noindent\textbf{Solution 9.8} (IOAA 2021 TH Q6): source uses 28.55 pc above but 28.57 here
\smallskip

\noindent\textbf{Solution 9.8} (IOAA 2021 TH Q6): source misprint for "minimum"
\smallskip

\noindent\textbf{Solution 9.8} (IOAA 2021 TH Q6): source gives a = 7.25e10 m above but uses 7.21e10 m here, and v\_p = 40 km/s above but 36000 m/s here
\smallskip

\noindent\textbf{Solution 9.10} (IOAA 2021 DA DA2): source spelling "Stephan-Boltzmann"
\smallskip

\noindent\textbf{Solution 9.10} (IOAA 2021 DA DA2): source wording "which is give by"
\smallskip

\noindent\textbf{Solution 9.10} (IOAA 2021 DA DA2): source writes "5778 $^\circ$K"
\smallskip

\noindent\textbf{Solution 9.10} (IOAA 2021 DA DA2): source wording
\smallskip

\section*{Chapter 10 --- Compact Objects \& Relativistic Astrophysics}

\noindent\textbf{Solution 10.9} (IOAA 2016 TH T10): "rquirement" in source
\smallskip

\noindent\textbf{Solution 10.14} (IOAA 2022 TH 13): the official solution writes sin(theta\_L) here; from part (c) the angle is theta = omega*t, not theta\_L.
\smallskip

\noindent\textbf{Solution 10.14} (IOAA 2022 TH 13): the official solution omits the factor omega inside the cosine; it should read cos[omega(t\_L - R/c)].
\smallskip

\noindent\textbf{Solution 10.14} (IOAA 2022 TH 13): the official solution substitutes 0.0975 = 1/gamma\^{}2 where the formula above calls for gamma = 3.2; evaluating [gamma(1-v/c)]\^{}{-4} directly gives \textasciitilde{}1.5e3, not 1.8e9 (and similarly for A\_min below).
\smallskip

\noindent\textbf{Solution 10.19} (IOAA 2025 TH T12): source writes dM\_bh/M\_bh\^{}2, though the integral that follows is of M\_bh\^{}2 dM\_bh
\smallskip

\section*{Chapter 11 --- Galactic Astrophysics}

\noindent\textbf{Solution 11.2} (IOAA 2009 TH P17): the exponent is printed with m+5.2 here but m-5.2 in the final expressions below; with x in kpc the correct form is 10\^{}{(m-5.2)/5}.
\smallskip

\noindent\textbf{Solution 11.8} (IOAA 2015 TH S10): the first term is printed as (2/5)m\_H\^{}2; it should read (2/5)m\_H r\_H\^{}2
\smallskip

\noindent\textbf{Solution 11.8} (IOAA 2015 TH S10): the source prints 1.38e-16 (CGS k) although the numerical result follows from the SI value 1.38e-23; "6,40" (decimal comma) as printed.
\smallskip

\noindent\textbf{Solution 11.10} (IOAA 2020 TH 5): the source prints "$\sigma$ 10\^{}{-21}" in the denominator
\smallskip

\noindent\textbf{Solution 11.13} (IOAA 2014 DA D1): the text prints (i = 0) here although it describes the general case of a nonzero inclination.
\smallskip

\noindent\textbf{Solution 11.13} (IOAA 2014 DA D1): comma
\smallskip

\noindent\textbf{Solution 11.13} (IOAA 2014 DA D1): the tables are captioned "Table 2 The coordinates in scaled plane"
\smallskip

\noindent\textbf{Solution 11.13} (IOAA 2014 DA D1): "maesured"/"mesured" as printed
\smallskip

\noindent\textbf{Solution 11.13} (IOAA 2014 DA D1): 0,01416 above vs 0,014666 here, as printed
\smallskip

\noindent\textbf{Solution 11.13} (IOAA 2014 DA D1): "43 m" as printed (should be mm)
\smallskip

\noindent\textbf{Solution 11.13} (IOAA 2014 DA D1): as printed (the usual value is 2.0e30 kg)
\smallskip

\noindent\textbf{Solution 11.15} (IOAA 2024 DA D2): decimal commas as printed
\smallskip

\section*{Chapter 12 --- Extragalactic Astrophysics}

\noindent\textbf{Solution 12.1} (IOAA 2008 TH S15): source says "pc" and "8.41e19 cm"; from the definition R25 is in kpc (27.2678 kpc = 8.41e22 cm), and the final mass below is consistent with kpc.
\smallskip

\noindent\textbf{Solution 12.1} (IOAA 2008 TH S15): mixed CGS/SI units as in the original
\smallskip

\noindent\textbf{Solution 12.2} (IOAA 2011 TH L2): source has "+"; the numerical substitution below (and the first line) require "-"
\smallskip

\noindent\textbf{Solution 12.3} (IOAA 2012 TH L2): as printed; the prefactor 1/2 and the denominator 80 do not reproduce 23.9 (1/2 x 11.54 = 5.77), which is closer to 2 x 11.54
\smallskip

\noindent\textbf{Solution 12.4} (IOAA 2015 TH S7): G value and units as in the original (should be 6.674e-11 N m\^{}2 kg\^{}-2)
\smallskip

\noindent\textbf{Solution 12.7} (IOAA 2018 TH T8): "viral" for "virial" as in the original
\smallskip

\noindent\textbf{Solution 12.7} (IOAA 2018 TH T8): "dym" as in the original
\smallskip

\noindent\textbf{Solution 12.8} (IOAA 2023 TH Q11): "bremmstrahlung" as in the original
\smallskip

\noindent\textbf{Solution 12.8} (IOAA 2023 TH Q11): "the the" as in the original
\smallskip

\noindent\textbf{Solution 12.8} (IOAA 2023 TH Q11): parts (c)+(d) point values in the source solution (4+3+2) do not match the statement's split (3+6); transcribed as printed.
\smallskip

\noindent\textbf{Solution 12.11} (IOAA 2011 DA D2): as printed; (290.8/2)/cos(35 deg) = 177.5 km/s, not 216.7
\smallskip

\noindent\textbf{Solution 12.12} (IOAA 2016 DA D3): 22.2 as printed, although Delta t\_0 is quoted as 22.0 d above
\smallskip

\section*{Chapter 13 --- Cosmology}

\noindent\textbf{Solution 13.9} (IOAA 2013 TH L1): 8.6e-27 kg/m\^{}3 is the critical density for H\_0 = 67.8; the matter density with the 0.32 factor would be 2.8e-27 kg/m\^{}3.
\smallskip

\noindent\textbf{Solution 13.16} (IOAA 2018 TH T11): source writes (ML\^{}2 T\^{}{-2})\^{}x for hbar despite stating [ML\^{}2 T\^{}{-1}] above; the exponent bookkeeping below uses T\^{}{-1}
\smallskip

\noindent\textbf{Solution 13.19} (IOAA 2021 TH Q15): "cosmic string trings" as in source
\smallskip

\noindent\textbf{Solution 13.20} (IOAA 2022 TH 4): source writes l(0) for the difference l(0) - l(t\_0)
\smallskip

\noindent\textbf{Solution 13.23} (IOAA 2025 TH T10): misplaced parenthesis in source; the numerical line below has the intended (t\_nuc - t\_wk)/880.6
\smallskip

\noindent\textbf{Solution 13.25} (IOAA 2017 DA D2): source table header says (kg); the values are in solar masses
\smallskip

\noindent\textbf{Solution 13.25} (IOAA 2017 DA D2): source labels y as M\_DM(r); the tabulated values are M\_DM(r)/r
\smallskip

\noindent\textbf{Solution 13.27} (IOAA 2022 DA 2): "We can phi values" as in source (missing "use")
\smallskip

\section*{Chapter 14 --- Gravitational Waves}

\noindent\textbf{Solution 14.1} (IOAA 2016 TH T11): the official solution prints the mass ratio below as M\_WD/M\_sun; from R $\propto$ M\^{}{-1/3} it should be M\_sun/M\_WD, which is what the subsequent lines (and the numerical answer) actually use.
\smallskip

\noindent\textbf{Solution 14.1} (IOAA 2016 TH T11): the official solution says 166.67 Hz here, although part (T11.1) found 142.86 Hz (and 142.86 is used in the mass estimate below).
\smallskip

\noindent\textbf{Solution 14.5} (IOAA 2022 DA 1): the official solution labels this second value m\_{2,min}; it is the minimum of the primary mass, m\_{1,min}.
\smallskip

\noindent\textbf{Solution 14.5} (IOAA 2022 DA 1): the official solution omits the minus sign; the upper value should be -7.8 s (= 2.2 s - 10 s).
\smallskip

\noindent\textbf{Solution 14.5} (IOAA 2022 DA 1): correspondingly, the upper value should be -1.8e-15.
\smallskip
